% Ad Familiares 1.1 (ad-familiares-01-01)
% Source: https://raw.githubusercontent.com/PerseusDL/canonical-latinLit/master/data/phi0474/phi056/phi0474.phi056.perseus-lat1.xml
% Fetched by scripts/fetch_latin.py

% Dateline (Perseus): Scr. Romae Id. Ian. a. u. c. 698 (a. Chr. 56) .

\ciceroLetterOpener{M. CICERO S. D. P. LENTVLO PROCOS.}

\ciceroSection{1} ego omni officio ac potius pietate erga te ceteris satis facio omnibus, mihi ipse numquam satis facio; tanta enim magnitudo est tuorum erga me meritorum, ut, quod tu nisi perfecta re de me non conquiesti, ego quia non idem in tua causa efficio, vitam mihi esse acerbam putem. in causa haec sunt: Hammonius, regis legatus, aperte pecunia nos oppugnat; res agitur per eosdem creditores, per quos, cum tu aderas, agebatur; regis causa si qui sunt qui velint, qui pauci sunt, omnes rem ad Pompeium deferri volunt, senatus religionis calumniam non religione, sed malevolentia et illius regiae largitionis invidia comprobat.

\ciceroSection{2} Pompeium et hortari et orare et iam liberius accusare et monere, ut magnam infamiam fugiat, non desistimus; sed plane nec precibus nostris nec admonitionibus relinquit locum, nam cum in sermone cotidiano tum in senatu palam sic egit causam tuam, ut neque eloquentia maiore quisquam nec gravitate nec studio nec contentione agere potuerit, cum summa testificatione tuorum in se officiorum et amoris erga te sui. Marcellinum tibi esse iratum scis; is hac regia causa excepta ceteris in rebus se acerrimum tui defensorem fore ostendit. quod dat, accipimus; quod instituit referre de religione et saepe iam retulit, ab eo deduci non potest.

\ciceroSection{3} res ante Idus acta sic est (nam haec Idibus mane scripsi): Hortensi et mea et Luculli sententia cedit religioni de exercitu; teneri enim res aliter non potest; sed ex illo senatus consulto, quod te referente factum est, tibi decernit, ut regem reducas, quod commodo rei p. facere possis, ut exercitum religio tollat, te auctorem senatus retineat. Crassus tris legatos decernit nec excludit Pompeium, censet enim etiam ex iis, qui cum imperio sint; Bibulus tris legatos ex iis, qui privati sunt. huic adsentiuntur reliqui consulares praeter Servilium, qui omnino reduci negat oportere, et Volcacium, qui Lupo referente Pompeio decernit, et Afranium, qui adsentitur Volcacio. quae res auget suspicionem Pompei voluntatis, nam advertebatur Pompei familiaris adsentiri Volcacio. laboratur vehementer; inclinata res est. Libonis et Hypsaei non obscura concursatio et contentio omniumque Pompei familiarium studium in eam opinionem rem adduxerunt, ut Pompeius cupere videatur; cui qui nolunt, idem tibi, quod eum ornasti, non sunt amici;

\ciceroSection{4} nos in causa auctoritatem eo minorem habemus, quod tibi debemus, gratiam autem nostram exstinguit hominum suspicio, quod Pompeio se gratificari putant. ut in rebus multo ante, quam profectus es, ab ipso rege et ab intimis ac domesticis Pompei clam exulceratis, deinde palam a consularibus exagitatis et in summam invidiam adductis ita versamur. nostram fidem omnes, amorem tui absentis praesentes tui cognoscent. si esset in iis fides, in quibus summa esse debebat, non laboraremus.
